\subsection{表示的系数}

设 U 是 K 上带单位元的结合代数，U* 是向量空间 U 的对偶，\(\rho\) 是 U 在向量空间 E 上的一个表示。对于 \(e \in E\) 和 \(e' \in E^*\)，令 \(\theta(e, e') \in U^*\) 为 \(\rho\) 的相应系数（《代数》，第八章，§ 13，第 3 号）。回顾，\(\theta(e, e')(x) = \langle \rho(x)e, e' \rangle\)；并且映射 \(e \mapsto \theta(e, e')\) 是在 \(e'\) 固定时，从 U-模 E 到具有 U 的余正则表示的 U-模 U* 的同态（同上，命题 1）；因此，由 \(\rho\) 的系数生成的 U* 向量子空间（本段记为 C(\(\rho\))）是 U* 的子 U-模。若 \((e_i')_{i \in I}\) 是在 K 上生成 E* 的一族元素，则映射 \(e \mapsto (\theta(e, e_i'))\) 是从 E 到 C(\(\rho\))\(^{I}\) 的单射 U-同态；因为对于所有 i，\(\theta(e, e_i') = 0\) 就意味着 \(\langle e, e_i' \rangle = \langle \rho(1)e, e_i' \rangle = \theta(e, e_i')(1) = 0\) 对所有 i 成立，从而 e=0。

特别地，若 U 是李代数 g 的包络代数，且 \(\rho\) 是 g 的表示（视为 U 的表示），作用在 n 维向量空间 E 上，则 g-模 E 同构于 \((\mathrm{C}(\rho))^{n}\) 的一个子 g-模。

